\chapter{Introducción}

Este manual explica cómo usar el sistema de gestión de \textbf{Brasa Brava}, el que se usa en el local para tomar los pedidos de cada mesa, mandarlos a cocina, cobrar, revisar las ventas y mantener al día el menú, el inventario y el personal. Está pensado para leerse de corrido la primera vez y para consultarse después por partes: cada capítulo corresponde a una pantalla del menú.

\section{Qué hace el sistema}

\begin{tabla}{>{\titulos\bfseries}l X}
  \encabezado{Pantalla} & \encabezado{Para qué sirve} \\
  Página principal & Saludo, hora del local, resumen de caja y cocina del día y accesos rápidos \\
  Familia & Catálogo del local: productos, recetas, porciones que alcanzan con el stock y promociones \\
  Caja & Plano de mesas, armado del pedido, comanda para cocina, lo que ya salió, cobro y ticket \\
  Reportes & Ventas cobradas por fechas y por empleado, en PDF o por correo, y envíos automáticos \\
  Productos & Menú: precio, foto, categoría, disponibilidad y receta de cada producto \\
  Secciones & Ambientes del local y sus mesas con su capacidad \\
  Inventario & Insumos con su cantidad y su mínimo \\
  Categorías & Categorías y subcategorías del menú \\
  Promociones & Combos a precio fijo con sus fechas, días y horario \\
  Empleados & Personal, sus datos, su cargo y su acceso al sistema \\
  Cargos & Qué pantallas ve cada cargo \\
  Eliminación física & Borrado definitivo de lo dado de baja (solo el DIRECTORIO) \\
\end{tabla}

\section{Escritorio y web}

El sistema funciona de dos formas con las mismas pantallas y los mismos datos de cada local:

\begin{description}
  \item[Aplicación de escritorio (Windows).] Se abre como cualquier programa desde su ícono. Trae su propia base de datos, así que \textbf{no necesita internet} para vender. Es la que imprime directo en las ticketeras de cocina y de caja, sin abrir el diálogo de impresión. Al actualizarla no se pierde ningún dato.
  \item[Versión web.] Se abre en el navegador con la dirección que entregue el encargado del sistema. Imprime con el diálogo de impresión del navegador, donde se elige la ticketera.
\end{description}

Las pantallas se adaptan a computadora, tablet y celular. En pantallas angostas el menú se abre desde la barra naranja y las tablas se muestran como tarjetas (sección~\ref{sec:celular}).

\section{Quién ve qué}

Cada empleado tiene un \textbf{cargo}, y cada cargo tiene marcadas las pantallas que puede ver (sección~\ref{sec:cargos}). Si una pantalla no aparece en el menú es porque el cargo no la tiene. El \textbf{DIRECTORIO} es la cuenta del dueño: ve todo, cambia su propia contraseña, programa los envíos automáticos de reportes y es el único que puede hacer la eliminación física.

\begin{tabla}{>{\titulos\bfseries}l X}
  \encabezado{Cargo} & \encabezado{Pantallas con las que se instala el sistema} \\
  Administrador & Todas, menos la eliminación física \\
  Cajero & Página principal, Familia y Caja \\
  Mesero & Página principal, Familia y Caja \\
  Cocinero & Página principal y Familia \\
  DIRECTORIO & Todas, incluida la eliminación física \\
\end{tabla}

Los cargos se pueden cambiar y se pueden crear otros (por ejemplo, Bartender o Delivery) en \pantalla{Administración \guilsinglright{} Cargos}.

\section{Cómo leer este manual}

\begin{itemize}
  \item Los botones se escriben así: \boton{Enviar a cocina}.
  \item Las pantallas y menús, así: \pantalla{Administración \guilsinglright{} Productos}.
  \item Los campos de un formulario, en negrita: \campo{Correo de destino}.
  \item Los montos están en bolivianos (Bs) con coma decimal: Bs 85,00.
  \item Las fechas se muestran como día/mes/año y las horas en formato de 12 horas (AM/PM), siempre en la hora del local.
  \item Las capturas se tomaron con los datos de demostración del sistema; los nombres y montos de su local serán otros.
\end{itemize}

\begin{consejo}
Los campos obligatorios llevan un asterisco (*). Si falta algo o está mal escrito, el formulario no se guarda y el mensaje aparece en rojo debajo del campo, diciendo qué corregir.
\end{consejo}
